\documentclass[crop=false]{standalone}
% --- ПРЕАМБУЛА ДЛЯ ПОДДЕРЖКИ РУССКОГО ЯЗЫКА И МАТЕМАТИКИ ---
\usepackage[T2A]{fontenc}
\usepackage[utf8]{inputenc}
\usepackage[russian]{babel}
\usepackage{amsmath}
\usepackage{geometry}
\usepackage{amssymb}
\usepackage{graphicx}
\usepackage{float}
\usepackage{import}
\geometry{a4paper, top=2cm, bottom=2cm, left=2cm, right=2cm}
\newcommand{\sidecomment}[1]{\marginpar{\raggedright\scriptsize\itshape #1}}
\newcommand{\iu}{\mathrm{i}}

\newcommand{\Def}{%
    \text{Def}%
    \hspace{-0.85em}% Сдвиг назад ровно на середину (подберите под шрифт)
    \raisebox{0.1ex}{/}% Черта
    \hspace{1em}% Возврат к концу слова + небольшой отступ
}

\renewcommand{\Re}{\operatorname{Re}}
\renewcommand{\Im}{\operatorname{Im}}

\ifstandalone
    \newcommand{\lecturebreak}{} % В отдельном файле лекции разрыв не нужен
\else
    \newcommand{\lecturebreak}{\clearpage} % В полной книге - начинаем с новой страницы
\fi

\newcounter{lecturenumber} % Создаем счетчик для нумерации лекций
\newcommand{\lecture}[2]{
    \lecturebreak
    \refstepcounter{lecturenumber} % Увеличиваем счетчик на 1
    \par\vspace{2em}\noindent % Отступ сверху и убираем абзацный отступ
    {\Large\bfseries % Делаем заголовок крупным и жирным
    Лекция \thelecturenumber: #1 % Выводим "Лекция N: Тема"
    \hfill % Растягиваем пробел, чтобы дата была справа
    \large\textit{#2}} % Выводим дату курсивом
    \par\vspace{1em} % Небольшой отступ снизу
    \addcontentsline{toc}{section}{Лекция \thelecturenumber: #1} % Добавляем в оглавление
    \par
}

\newcommand{\TwoCols}[2]{%
  \par\noindent % Начать с новой строки без отступа
  \begin{minipage}[t]{0.48\textwidth}
    #1
  \end{minipage}% <-- Важный '%' для удаления пробела
  \hfill % Растягивающийся пробел между колонками
  \begin{minipage}[t]{0.48\textwidth}
    #2
  \end{minipage}% <-- Важный '%' для удаления пробела
  \par%
}

\pagestyle{plain} % Включаем нумерацию страниц\
\begin{document}
\lecture{}{05.06}
\Def Лагранжевы переменные
\[q_1, \dots, q_n, \dot q_1, \dots, \dot q_n\]
\Def Гамильтоновы переменные
\[q_1, \dots, q_n, \dot p_1, \dots, \dot p_n\]
\section{Интерпретация уравнений Гамильтона. Скобка Пуассона}
\Def Первый интеграл уравнения Гамильтона — непрерывно дифференцируемая функция, обращающаяся в константу при подстановке решения уравнения Гамильтона
\[f(p(t), q(t), t) = \text{const}\]

\noindent \keypoint 

$f$ — первый интеграл уравнения Гамильтона 

$\iff$
\begin{flalign*}
\frac{d f}{d t} =  \{f, H\} + \frac{\partial f}{\partial t} = 0 &&
\end{flalign*}

\begin{proof}
    \[
        \frac{d f}{d t} = \sum_{i=1}^{n} \frac{\partial f}{\partial q_i} \dot q_i + \sum_{i=1}^{n} \frac{\partial f}{\partial p_i} \dot p_i + \frac{\partial f}{\partial t} = \underbracket{\sum_{i=1}^{n} \left( \frac{\partial f}{\partial q_i }\frac{\partial H}{\partial p_i} - \frac{\partial f}{\partial p_i} \frac{\partial H}{\partial q_i}\right)}_{\{f, H\}} + \frac{\partial f}{\partial t}
        \]
    В обратную сторону по определению
\end{proof}
Замечание: Тождество из доказательства справедливо для любой функции Гамильтоновых переменных
\Def Скобка Пуассона $\{f, g\}$, где $f, g$ — функции Гамильтоновых переменных
% fixme поправить наложения
\[\{f, g\} = \sum_{i=1}^{n} \begin{vmatrix} \frac{\partial f}{\partial q_i} & \frac{\partial g}{\partial q_i} \\ \frac{\partial f}{\partial p_i} & \frac{\partial g}{\partial p_i} \end{vmatrix}\]
\keypoint Уравнения Гамильтона через скобки Пуассона
\[
    \dot q_i = \{q_i, H\} \qquad
    \dot p_i = \{p_i, H\}
    \]
\begin{proof}
    \[\{q_i, H\} = \sum_{j=1}^{n} \begin{vmatrix}
        \frac{\partial q_i}{\partial q_j} & \frac{\partial H}{\partial q_j} \\
        \frac{\partial q_i}{\partial p_j} & \frac{\partial H}{\partial p_j}  
    \end{vmatrix} = \frac{\partial H}{\partial p_i} \qquad 
    \{p_i, H\} = \sum_{j=1}^{n} \begin{vmatrix}
        \frac{\partial p_i}{\partial q_j} & \frac{\partial H}{\partial q_j} \\
        \frac{\partial p_i}{\partial p_j} & \frac{\partial H}{\partial p_j}
    \end{vmatrix} = - \frac{\partial H}{\partial q_i}\]
\end{proof}

\noindent \keypoint Свойства скобки Пуассона
\begin{enumerate}
    \item $\{f, g\} = - \{g, f\}$
    \item $\{c f, g\} = \{ f, c g \} = c \{f , g\}$
    \item $\frac{\partial}{\partial t} \{f, g\} = \{\frac{\partial f}{\partial t}, g\} + \{f, \frac{\partial g}{\partial t}\}$
    \item $\{\{f, g\}, h\} + \{\{g, h\}, f\} + \{\{h, f\}, g\} = 0$
\end{enumerate}
\begin{proof}
    1, 2, 3 очевидно
    \begin{enumerate}
        \setcounter{enumi}{3} 
        \item Заметим:
        \[
        \begin{aligned}
            \{\{f, g\}, h\} &= \{f, g\}_q h_p - \{f, g\}_p h_q = \left(f_q g_p - f_p g_q \right)_q h_p - \left(f_q g_p - f_p g_q \right)_p h_q \\
            &= \left(f_{qq} g_p + f_q g_{qp} - f_{qp} g_{q} - f_p g_{qq}\right) h_p - \left(f_{pq} g_p + f_q g_{pp} - f_{pp} g_q - f_p g_{pq} \right) h_q
        \end{aligned}\]
        \[
            \{\{g, h\}, f\} = \left(g_{qq} h_p + g_q h_{qp} - g_{qp} h_{q} - g_p h_{qq}\right) f_p - \left(g_{pq} h_p + g_q h_{pp} - g_{pp} h_q - g_p h_{pq} \right) f_q
        \]
        \[
            \{\{h, f\}, g\} = \left(h_{qq} f_p + h_q f_{qp} - h_{qp} f_{q} - h_p f_{qq}\right) g_p - \left(h_{pq} f_p + h_q f_{pp} - h_{pp} f_q - h_p f_{pq} \right) g_q
        \]
    \end{enumerate}
\end{proof}

\begin{theorem}[Якоби-Пуассона]
    Если $f,g$ — первые интегралы уравнения Гамильтона
    
    То  $\{f, g\}$ — первый интеграл уравнения Гамильтона
\end{theorem}
\begin{proof}
    \[\{f, H\} + \frac{\partial f}{\partial t} = 0 \qquad \{g, H\} + \frac{\partial g}{\partial t} = 0\]
    Автору конспекта неудобно писать эти скобки, далее $\{\{a, b\}, c\} = \{a, b , c\}$
    \[\{\{f, g\}, H\} + \frac{\partial \{f, g\}}{\partial t} = \{f, g, H\} + \{\frac{\partial f}{\partial t}, g\} + \{f, \frac{\partial g}{\partial t}\}  = \{f, g, H\} + \{-\{f , H\}, g\} + \{f, -\{g, H\}\} = 0\]
\end{proof}
Замечание:
Достаточно найти $n$ первых интегралов, остальные выражаются через них

Замечание:
Теорема Я-П полезна в численных методах, там накапливаются ошибки, полезно следить за первым интегралом для оценки ошибки

\section{Интеграл Якоби}
\noindent \keypoint Цитата Сергея Алексеевича: Ньютон кстати может и не математик в первую очередь, он же глава королевского монетного двора, может он вообще в первую очередь был политиком, вы знали что Ньютон был трейдером: купил кучу акций Ост-Индской компании и потерял кучу денег

\noindent \keypoint
\[H \text{ — I инт} \iff \frac{\partial H}{\partial t} = 0\]
\begin{proof}
    \[\frac{d}{d t} H = \underbracket{\{H, H\}}_0 + \frac{\partial H}{\partial t} \qquad \frac{d H}{d t} = \frac{\partial H}{\partial t} \qquad \]
\end{proof}
\noindent \keypoint В стационарных системах Гамильтониан постоянен во времени и \emph{в некотором смысле} является полной энергией

\noindent \keypoint

для стационарных систем 

Если $f$ — первый интеграл уравнения Гамильтона

То

\[\frac{\partial^2 f}{\partial t^2}  = \{H, \{H, f\}\}\]
\begin{proof}
    \[\frac{\partial H}{\partial t} := 0, f \text{ — I инт} \implies \frac{\partial f}{\partial t} =-\{f, H\} = \{H, f\}\]
\end{proof}
Замечание: В стационарных системах любая частная производная первого интеграла по времени является первым интегралом
\section{Циклические координаты}
\Def Обобщённая координата называется циклической, если Лагранжиан(т.е. и Гамильтониан) от неё не зависит явно

\noindent \keypoint Обобщённые импульсы соответсвующие циклическим координатам являются первыми интегралами системы
\begin{proof}
\[\frac{d}{d t} \underbracket{\frac{\partial L}{\partial \dot q_i}}_{p_i} + \underbracket{\frac{\partial L}{\partial q_i}}_0 \implies \dot p_i = 0\]
Из предыдущих утверждений следует, что $p_i$ — первый интеграл
\end{proof}
\noindent \keypoint Цитата Сергея Алексеевича: за вас никто учиться не будет, нейросети ..., ну кстати нейросети могут за вас учиться

\section{Динамические уравнения Эйлера}

Динамические уравнения эйлера для твёрдого тела с неподвижной точкой
\[
\begin{cases}
    A \dot p + q r (C - B) = M_{O x} \\
    B \dot q + p r (A - C)= M_{O y} \\
    C \dot r + p q ( B - A) = M_{O z}
\end{cases}
\]
\begin{proof}

$K', M'$ — кинетический момент и главный момент внешних сил, записанные в неподвижной системе координат

$K, M$ — записанные в системе, связанной с твёрдым телом
\[\bar K_O' = A \bar K_O \qquad \bar M_O' = A \bar M_O \qquad \frac{d}{d t} A \bar K_O = A \bar M_O \]
\[A \dot {\bar K}_O + \bar \omega' \times A \bar K_O = A \bar M_O \big| A^T \cdot\]
\[\dot {\bar K}_O + \bar \omega \times \bar K_O = \bar M_O \qquad \dot {\bar K}_O =  J \dot{\bar \omega} = \begin{pmatrix}
    A \dot p \\ B \dot q \\ C \dot r
\end{pmatrix}\]
\end{proof}
Замечание: Можно взять решения Динамических и подставить в кинематические, чтобы найти движение тела, 

можно подставить решения кинематических в динамические, чтобы исследовать устойчивость вращения тела

Случай Эйлера: на систему не действуют внешние силы, то есть правая часть равна нулю

\section{Первые интегралы в случае Эйлера}
\noindent \keypoint Кинетическая энергия — первый интеграл
\[T = \frac{1}{2} \left(A p^2 + B q^2 + C r^2 \right)= \text{const}\]
\begin{proof}
    Из теоремы об изменении кинетической энергии
\end{proof}
\noindent \keypoint Модуль кинетического момента относительно неподвижной точки твёрдого тела — первый интеграл
\[\left| \begin{pmatrix}
    A p \\ B q \\ C r
\end{pmatrix}\right| = \text{const} \qquad A^2 p^2 + B^2 q^2 + C^2 r^2 = \text{const}\]
\begin{proof}
    Из теоремы об изменении кинетического момента
\end{proof}

Замечание: Соотношения между $A, B, C$ характеризуют различные эффекты, например эффект Джанибекова

\end{document}